\documentclass[11pt]{article}
\usepackage[margin=1in]{geometry}
\usepackage{amsmath,amssymb,amsthm}
\usepackage{graphicx}
\usepackage{xcolor}
\usepackage{microtype}
\usepackage{enumitem}
\usepackage[colorlinks=true,linkcolor=blue!55!black,citecolor=blue!55!black,urlcolor=blue!55!black]{hyperref}

\newtheorem{theorem}{Theorem}
\newtheorem{lemma}{Lemma}
\newcommand{\E}{\mathbb E}
\newcommand{\R}{\mathbb R}
\newcommand{\muinf}{\mu_\infty}

\title{\textbf{Sharp Fourier restriction exponent for Mandelbrot cascades}\\
\large The squared-critical and squared-supercritical phases}
\author{Automated mathematical discovery run \texttt{fa\_banach\_001}\\
Lane 7, model GPT5.6}
\date{23 July 2026}

\begin{document}
\maketitle

\begin{center}
\fbox{\begin{minipage}{0.91\textwidth}
\textbf{Status.} Candidate partial result, likely valid, pending expert review.
The result determines the supremal restriction exponent in the
squared-critical and squared-supercritical regimes.  It does not decide
endpoint attainment or the remaining squared-subcritical regime.
\end{minipage}}
\end{center}

\section{Source question}

The source is Xinxin Chen, Yong Han, Yanqi Qiu, and Zipeng Wang,
\emph{Harmonic analysis of Mandelbrot cascades -- in the context of
vector-valued martingales}, arXiv:2409.13164v3 \cite{CHQW}.
After Corollary 1.6, on page 9, the authors ask:
\begin{quote}
What is the optimal exponent \(r\in(1,2)\) in
\[
 \|\widehat f\|_{L^2(\muinf)}
 \leq C(r,\muinf)\|f\|_{L^r(\R)}?
\]
\end{quote}

\begin{figure}[ht]
\centering
\includegraphics[width=\textwidth]{figures/open_problem_crop.png}
\caption{Corollary 1.6 and the open question in arXiv:2409.13164v3, page 9.}
\end{figure}

\section{Definitions and assumptions}

Fix an integer \(b\geq2\).  Let \(W\geq0\) be nonconstant with
\(\E W=1\), and let \(\muinf\) be the associated canonical \(b\)-adic
Mandelbrot cascade on \([0,1]\).  We use the source assumptions
\[
 \E[W\log W]<\log b,
 \qquad
 \E[W^t]<\infty\quad(t>0),
\]
and, in the squared-supercritical case, the source's nonlattice assumption
on \(\log W\).  All almost-sure statements below are on
\(\{\muinf\neq0\}\).

Set
\[
 W^{(2)}=\frac{W^2}{\E W^2}.
\]
The squared-critical and squared-supercritical regimes are characterized by
\[
 \E[W^{(2)}\log W^{(2)}]\geq\log b.
\]

For a finite Borel measure \(\mu\) on \(\R\), define its Frostman dimension by
\[
 \dim_{\mathrm{Fr}}\mu
 =\sup\left\{\alpha\geq0:
 \mu(I)\leq C_\alpha |I|^\alpha
 \text{ for every interval }I\right\}.
\]
For \(1\leq r<2\), let \(\mathcal R_r(\mu)\) denote the restriction estimate
\[
 \|\widehat f\|_{L^2(\mu)}\leq C_r\|f\|_{L^r(\R)}
 \qquad(f\in\mathcal S(\R)).
\]
Define the supremal restriction exponent
\[
 r_*(\mu)=\sup\{r\in[1,2):\mathcal R_r(\mu)\text{ holds}\}.
\]

\section{Idea of the proof}

The source already proves the lower half of the answer:
\(\mathcal R_r(\muinf)\) holds for
\(r<4/(4-D_F)\).  The missing sharpness follows from local mass.
A restriction estimate at exponent \(r\) forces every interval of length
\(\delta\) to have mass \(O(\delta^{2/r'})\).  Thus restriction cannot pass
the exact Frostman dimension.

For Mandelbrot cascades, the Frostman exponent is the maximum of a simple
moment profile.  The squared-critical condition says that this profile has
already stopped increasing by \(q=2\).  Rewriting the source's
squared-supercritical formula then shows that the Fourier dimension is
exactly twice the Frostman dimension.  The localization obstruction and the
source's positive estimate consequently meet at the same exponent.

\section{The partial result}

\begin{theorem}[Sharp supremal exponent]\label{thm:main}
Under the assumptions above, suppose
\[
 \E[W^{(2)}\log W^{(2)}]\geq\log b.
\]
Let \(D_F=\dim_F(\muinf)\), whose deterministic value is given by
Theorem 1.1 of \cite{CHQW}.  Then almost surely on nonextinction,
\[
 r_*(\muinf)=\frac{4}{4-D_F}.
\]
More precisely, \(\mathcal R_r(\muinf)\) holds for every
\(1\leq r<4/(4-D_F)\), and it fails for every
\(r>4/(4-D_F)\).
\end{theorem}

\begin{proof}
Write
\[
 L(q)=\log \E W^q,\qquad
 a(q)=\frac{q-1-\log_b\E W^q}{q}
 \quad(q>1).
\]
The classical multifractal formula for a nondegenerate Mandelbrot cascade
gives, almost surely on nonextinction,
\begin{equation}\label{eq:frost}
 \dim_{\mathrm{Fr}}\muinf=\sup_{q>1}a(q).
\end{equation}
This is the extremal local-dimension, or uniform Holder, part of the standard
cascade spectrum; see Molchan \cite{Molchan} and the sharp modulus results of
Barral--Kupiainen--Nikula--Saksman--Webb \cite{BKNSW}.

We first locate the maximum in \eqref{eq:frost}.  Direct differentiation gives
\[
 a'(q)=
 \frac{\log b+L(q)-qL'(q)}{q^2\log b}.
\]
The numerator
\[
 H(q)=\log b+L(q)-qL'(q)
\]
is nonincreasing, since \(H'(q)=-qL''(q)\leq0\).  Nondegeneracy gives
\[
 H(1)=\log b-\E[W\log W]>0.
\]
On the other hand,
\[
 \E[W^{(2)}\log W^{(2)}]
 =2L'(2)-L(2)\geq\log b
\]
is exactly \(H(2)\leq0\).  Hence \(a\) attains its global maximum at some
\(q_*\in(1,2]\), and therefore
\begin{equation}\label{eq:max12}
 \sup_{q>1}a(q)=\sup_{1<q\leq2}a(q).
\end{equation}

We next compare this maximum with the source formula for \(D_F\).  In the
squared-supercritical case, formula (1.4) of \cite{CHQW} is
\[
 D_F
 =1-\inf_{1/2\leq t\leq1}
 \frac{\log\E[b^{1-t}W^{2t}]}{t\log b}.
\]
Putting \(q=2t\) shows that the expression being maximized is
\[
 1-\frac{\log\E[b^{1-t}W^{2t}]}{t\log b}
 =2a(q).
\]
Thus \eqref{eq:frost} and \eqref{eq:max12} give
\begin{equation}\label{eq:key}
 D_F=2\dim_{\mathrm{Fr}}\muinf.
\end{equation}
In the squared-critical case, \(H(2)=0\), so the maximum is at \(q=2\), and
the first branch of formula (1.4) yields the same identity:
\[
 D_F=1-\log_b\E W^2=2a(2)
 =2\dim_{\mathrm{Fr}}\muinf.
\]

It remains to convert \eqref{eq:key} into restriction sharpness.  Lemma
\ref{lem:localization} below shows that \(\mathcal R_r(\muinf)\) implies
\[
 \dim_{\mathrm{Fr}}\muinf\geq\frac{2}{r'}.
\]
Using \eqref{eq:key},
\[
 \frac{2}{r'}\leq\frac{D_F}{2},
\]
which is equivalent to
\[
 r\leq\frac{4}{4-D_F}.
\]
Therefore the restriction estimate fails for every larger \(r\).
Conversely, Corollary 1.6 of \cite{CHQW} proves it for every strictly smaller
\(r\).  This proves the asserted value of the supremum.
\end{proof}

\begin{lemma}[Frequency localization]\label{lem:localization}
Let \(\mu\) be a finite Borel measure on \(\R\).  If
\(\mathcal R_r(\mu)\) holds for some \(1\leq r<2\), then
\[
 \mu(I)\lesssim |I|^{2/r'}
\]
for every interval \(I\) of length at most \(1\).  In particular,
\(\dim_{\mathrm{Fr}}\mu\geq2/r'\).
\end{lemma}

\begin{proof}
Choose a Schwartz function \(\phi\) such that
\(\widehat\phi\geq1\) on \([-1,1]\).  If \(I\) has center \(\xi_0\) and
length \(\delta\leq1\), put
\[
 f_{\delta,\xi_0}(x)
 =\delta e^{2\pi i\xi_0x}\phi(\delta x).
\]
Then
\[
 \widehat f_{\delta,\xi_0}(\xi)
 =\widehat\phi\left(\frac{\xi-\xi_0}{\delta}\right),
 \qquad
 \|f_{\delta,\xi_0}\|_r
 =\delta^{1-1/r}\|\phi\|_r
 =\delta^{1/r'}\|\phi\|_r.
\]
Since \(\widehat f_{\delta,\xi_0}\geq1\) on \(I\), the assumed restriction
estimate gives
\[
 \mu(I)^{1/2}
 \leq\|\widehat f_{\delta,\xi_0}\|_{L^2(\mu)}
 \lesssim\delta^{1/r'}.
\]
Squaring proves the claim.
\end{proof}

\section{Explicit lognormal consequence}

For
\[
 W=e^{\sigma N-\sigma^2/2},\qquad N\sim\mathcal N(0,1),
\]
the nondegenerate range is \(0<\sigma^2<2\log b\), and the
squared-critical/supercritical range is
\(\sigma^2\geq\tfrac12\log b\).  In this range \cite{CHQW} gives
\[
 D_F
 =2\left(1-\frac{\sigma}{\sqrt{2\log b}}\right)^2.
\]
Consequently
\[
 r_*
 =\frac{4}{
 4-2\left(1-\frac{\sigma}{\sqrt{2\log b}}\right)^2}.
\]

\section{Verification notes and limitations}

\begin{enumerate}[leftmargin=*]
\item The localization lemma is deterministic and uses only Fourier scaling.
It also checks the direction of the necessary inequality: a larger \(r\)
forces a larger Frostman exponent.
\item The calculus step uses only convexity of
\(q\mapsto\log\E W^q\).  The squared phase condition is exactly the sign
condition \(H(2)\leq0\).
\item The identity \(D_F=2\dim_{\mathrm{Fr}}\muinf\) is special to the
squared-critical/supercritical phase.  In the squared-subcritical phase the
Frostman exponent can be strictly larger than \(D_F/2\), so the same argument
does not determine the optimal restriction exponent.
\item The proof identifies the supremum of valid exponents.  It does not
prove or disprove \(\mathcal R_{r_*}(\muinf)\) at the endpoint.
\item The exact Frostman spectrum formula \eqref{eq:frost} is an external
standard cascade theorem.  A reviewer should check its hypotheses against
the stated all-moments and nonlattice assumptions; these are stronger than
the usual hypotheses in the present regime.
\end{enumerate}

\section{Novelty check}

The bounded search covered the run's local 2000--2026 full-source index,
exact-phrase searches for the source question, searches combining
``Mandelbrot cascade'', ``optimal exponent'', ``restriction'', and
``supercritical'', and citations of arXiv:2409.13164 through 23 July 2026.
No paper stating Theorem \ref{thm:main} was found.

Carnovale--Fraser--de Orellana \cite{CFdO} give later \(L^q\)-dimension based
sufficient estimates and explicitly improve the original Stein--Tomas range
for some cascade parameters.  They do not state the matching localization
obstruction or determine the true critical exponent in the present phase.
Thus the result here is a direct combination of known exact dimensions with
a new sharpness identification, not a claimed new estimate in the
squared-subcritical regime.

\section{Human-review recommendation}

Promote as a substantial partial solution if the reviewer confirms the
standard exact Frostman formula \eqref{eq:frost} under the source hypotheses.
The most important review points are the convention for the cascade
\(L^q\)-spectrum, the equality case at squared criticality, and the
interpretation of ``optimal exponent'' as the supremum when endpoint
attainment is unknown.

\begin{thebibliography}{9}

\bibitem{CHQW}
X. Chen, Y. Han, Y. Qiu, and Z. Wang,
\emph{Harmonic analysis of Mandelbrot cascades -- in the context of
vector-valued martingales},
arXiv:2409.13164v3 (2025).

\bibitem{Molchan}
G. M. Molchan,
\emph{Scaling exponents and multifractal dimensions for independent random
cascades},
Commun. Math. Phys. \textbf{179} (1996), 681--702.

\bibitem{BKNSW}
J. Barral, A. Kupiainen, M. Nikula, E. Saksman, and C. Webb,
\emph{Critical Mandelbrot cascades},
Commun. Math. Phys. \textbf{325} (2014), 685--711;
arXiv:1206.5444.

\bibitem{CFdO}
M. Carnovale, J. M. Fraser, and A. E. de Orellana,
\emph{Fourier restriction estimates based on \(L^q\)-dimensions: beyond
Stein--Tomas},
arXiv:2606.07143 (2026).

\end{thebibliography}

\end{document}
